\documentclass[portuguese]{sbcreviews-2025}%
\usepackage[misc,geometry]{ifsym}
\usepackage{aas_macros}
\usepackage[bottom]{footmisc}
\usepackage{tabularray}
\usepackage{url}
\usepackage{pifont}
\usepackage{orcidlink} % Adicionado para suportar \orcidlink
\usepackage{xcolor}    % Adicionado para suportar \textcolor e row{1}={bg=gray!20}

\definecolor{engtitle}{RGB}{0,0,0}

\setcitestyle{square}

% --- Configurações do Periódico ---
\jid{ROCS}
\jtitle{SBC Computing Reviews, 2026, XX:1}
\issn{2966-3938}
\jyear{2026}

\category{Engenharia de Software / Inteligência Artificial}
\title[Rede Social Maker com Governança IoT e IA Generativa]{MakerConnect: Rede Social para Governança de Projetos IoT com Documentação Automatizada via Agentes de IA e n8n}
\engtitle{\textcolor{engtitle}{MakerConnect: A Maker Social Network for IoT Project Governance with Automated Documentation via AI Agents and n8n}}

\author[Fróes 2026]{
\affil{\textbf{Vinícius Fróes}~\orcidlink{0000-0000-0000-0000}~\textcolor{blue}{\faEnvelope}~~[{Católica de Santa Catarina}~|\href{mailto:vinicius.froes@catolicasc.edu.br}{~{\textit{vinicius.froes@catolicasc.edu.br}}}~]}
}

\begin{document}

\begin{frontmatter}
\maketitle

\begin{abstract-pt}
Este trabalho apresenta a MakerConnect, uma rede social técnica para a cultura Maker e Internet das Coisas (IoT) que integra um pipeline funcional de IA Generativa orquestrado via n8n. A plataforma aplica \textit{Retrieval-Augmented Generation} (RAG) com agentes baseados em GPT-4o para extrair requisitos de entradas não estruturadas e gerar automaticamente documentação técnica em PDF. A base de conhecimento vetorial (Pinecone/Supabase) é populada com dados reais de componentes eletrônicos, eliminando alucinações técnicas. Mecanismos sociais — \textit{forks}, \textit{upvotes} e logs de dificuldades — promovem reprodutibilidade e rastreabilidade. A conformidade com a LGPD é garantida por anonimização de PII antes do processamento pela IA. Resultados preliminares indicam redução expressiva do tempo de documentação e ampliação do reuso de hardware em projetos IoT.
\end{abstract-pt}

\begin{abstract-en}
This work presents MakerConnect, a technical social network for Maker culture and IoT that integrates a functional Generative AI pipeline orchestrated via n8n. The platform applies Retrieval-Augmented Generation (RAG) with GPT-4o-based agents to extract requirements from unstructured inputs and automatically generate PDF technical documentation. The vector knowledge base (Pinecone/Supabase) is populated with real electronic component data, preventing technical hallucinations. Social mechanisms — forks, upvotes, and error logs — promote reproducibility and traceability. LGPD compliance is ensured by anonymizing PII before AI processing. Preliminary results indicate significant reduction in documentation time and increased hardware reuse in IoT projects.
\end{abstract-en}

\begin{pchaves}
Cultura Maker, IA Generativa, RAG, n8n, Engenharia de Software, IoT, LGPD.
\end{pchaves}
\begin{keywords}
Maker Culture, Generative AI, RAG, n8n, Software Engineering, IoT, LGPD.
\end{keywords}

\end{frontmatter}

% ============================================================
\section{Introdução e Problema de Pesquisa}
% ============================================================

O movimento Maker enfrenta o fenômeno da \textit{documentação fantasma}: projetos IoT inovadores se perdem pela ausência de registros técnicos estruturados. Plataformas existentes (GitHub, Instructables) oferecem compartilhamento, mas não automatizam a produção documental nem orientam ativamente o maker. Isso resulta em baixa reprodutibilidade e reuso de hardware limitado. O problema de pesquisa é: \textit{como agentes de IA Generativa, orquestrados via n8n, podem reduzir essa lacuna documental, promovendo reprodutibilidade e rastreabilidade em comunidades maker?}

\textbf{Hipótese:} a integração de um pipeline RAG em uma rede social maker reduz significativamente o tempo e o esforço cognitivo da documentação técnica, aumentando padronização, reuso de componentes e reprodutibilidade de projetos IoT.

% ============================================================
\section{Objetivos}
% ============================================================

\textbf{Geral:} Desenvolver a MakerConnect, integrando agentes de IA Generativa e orquestração n8n para automatizar a documentação técnica e gerenciar o ciclo de vida de projetos IoT.

\textbf{Específicos:} (i)~Implementar pipeline funcional de IA com estágios de Extração, RAG e Pós-processamento; (ii)~Orquestrar workflows via n8n em nuvem (Hostinger); (iii)~Construir base de conhecimento vetorial com dados reais de componentes eletrônicos; (iv)~Modelar estrutura social com \textit{forks}, \textit{upvotes} e BOM via MySQL; (v)~Garantir conformidade LGPD com anonimização de PII; (vi)~Validar redução do esforço documental por métricas comparativas.

% ============================================================
\section{Abordagem de IA e Pipeline Funcional}
% ============================================================

A solução aplica \textbf{IA Generativa com RAG} — técnica da categoria \textit{Transformer/NLP Generativa} — integrada funcionalmente à plataforma, atendendo ao requisito de aplicação fundamentada em problema real. O pipeline segue quatro estágios obrigatórios, conforme a Tabela~\ref{tab:pipeline}.

\begin{table}[!ht]
\caption{Pipeline Funcional de IA da MakerConnect.}
\centering
\begin{tblr}{%
colspec={Q[l,0.18\linewidth] Q[l,0.36\linewidth] Q[l,0.36\linewidth]},
row{1}={bg=gray!20, font=\bfseries},
}
\hline
Estágio & Técnica / Ferramenta & Saída \\
\hline
Extração & Upload via React; NLP para parsing de descrições e imagens & Texto normalizado + metadados \\
Pré-proc. & Geração de \textit{embeddings} (GPT-4o); similaridade coseno; filtragem PII (LGPD) & Vetores indexados no Pinecone \\
Modelo IA & RAG: recuperação vetorial + geração com GPT-4o / Llama~3; prompt engineering com contexto técnico & Requisitos SW/HW estruturados \\
Pós-proc. & Renderização PDF; métricas de cobertura documental; logs de validação & Documentação auditável \\
\hline
\end{tblr}
\label{tab:pipeline}
\end{table}

\textbf{Justificativa do modelo:} RAG foi escolhido por combinar recuperação baseada em evidências com geração contextualizada, reduzindo alucinações técnicas críticas em domínios como eletrônica e IoT — onde erros de componente podem comprometer protótipos físicos. O uso de base vetorial própria (não dataset genérico) com dados reais de componentes garante aderência ao domínio.

\textbf{Base de dados:} a base de conhecimento é composta por datasheets e especificações reais de componentes eletrônicos (Arduino, ESP32, sensores, atuadores), coletados de fontes abertas e fabricantes. Dados são pré-processados, vetorizados e indexados no Pinecone. Nenhum dado de usuário é usado para ajuste do modelo — apenas metadados técnicos anonimizados.

% ============================================================
\section{Arquitetura da Solução}
% ============================================================

A MakerConnect adota arquitetura de \textbf{microserviços híbrida} com os seguintes componentes: \textit{frontend} em React (interface web para interação com a IA); backend em Node.js; n8n como orquestrador central de workflows em nuvem (Hostinger); banco relacional MySQL para dados sociais; e Pinecone/Supabase como Vector DB. A Tabela~\ref{tab:arch} resume as especificações técnicas.

\begin{table}[!ht]
\caption{Componentes Técnicos e Infraestrutura de IA.}
\centering
\begin{tblr}{%
colspec={Q[c,0.22\linewidth] Q[c,0.28\linewidth] Q[c,0.4\linewidth]},
row{1}={bg=gray!20, font=\bfseries},
}
\hline
Componente & Tecnologia & Função \\
\hline
Interface & React (Web) & Interação do maker com o pipeline de IA \\
Orquestrador & n8n / Hostinger & Coordenação de agentes e integração de APIs \\
Modelo LLM & GPT-4o + Llama~3 & Extração de requisitos e geração documental \\
Vector DB & Pinecone / Supabase & Base de conhecimento técnico (RAG) \\
Social DB & MySQL & Forks, upvotes, BOM, logs \\
Segurança & LGPD Middleware & Anonimização de PII pré-processamento \\
\hline
\end{tblr}
\label{tab:arch}
\end{table}

% ============================================================
\section{Trabalhos Relacionados}
% ============================================================

A presente pesquisa situa-se na interseção de três domínios emergentes: (i)~redes sociais técnicas para comunidades maker; (ii)~orquestração de agentes de IA generativa com \textit{Retrieval-Augmented Generation} (RAG); (iii)~governança e rastreabilidade de projetos IoT com conformidade regulatória (LGPD).

Para validar o posicionamento da MakerConnect, realizou-se busca sistemática em bases acadêmicas, retornando 3 trabalhos principais, publicados entre 2025 e 2026, que cobrem os pilares da proposta: aplicações de agentes assistentes em IoT (\cite{dong2025chatiot}); RAG federado para IoT restrito com foco em privacidade (\cite{hangyu2026federated}); e arquiteturas agentic emergentes para orquestração inteligente (\cite{singh2025agentic}). Plataformas de referência incluem Manual Maker (desafios maker com 1.5k+ participantes) e Rabbit Agents (orquestração IA com agentes automáticos).

A Tabela~\ref{tab:comparative} sintetiza o comparativo entre os 3 trabalhos acadêmicos de maior relevância, 2 plataformas comerciais e a MakerConnect, considerando dimensões-chave: foco em RAG, aplicação em IoT, orquestração, automação de documentação, camada social colaborativa, rastreabilidade técnica (fork/lineage), conformidade LGPD e status de implementação.

\begin{table}[!ht]
\caption{Matriz Comparativa: Trabalhos Relacionados, Plataformas e MakerConnect.}
\centering
\footnotesize
\begin{tblr}{%
width=\linewidth,
colspec={Q[l,1.95] Q[c,0.55] Q[c,0.55] Q[c,0.62] Q[c,0.55] Q[c,0.62] Q[c,0.62] Q[c,0.62] Q[c,0.92]},
row{1}={bg=gray!25, font=\bfseries},
row{2-4}={bg=blue!4},
row{5-6}={bg=orange!6},
row{7}={bg=green!12, font=\bfseries},
rowsep=1.5pt,
colsep=2.5pt,
}
\hline
Sistema/Trabalho & RAG & IoT & Orch. & Doc. & Social & Trace & LGPD & Status \\
\hline
Dong et al. (2025) & \ding{51} & \ding{51} & \ding{55} & \ding{51} & \ding{55} & \ding{55} & \ding{55} & Prot. \\
Hangyu (2026) & \ding{51} & \ding{51} & \ding{55} & \ding{55} & \ding{55} & \ding{55} & \ding{51} & Conc. \\
Singh (2025) & \ding{51} & \ding{55} & \ding{51} & \ding{55} & \ding{55} & \ding{55} & \ding{55} & Conc. \\
\hline
Manual Maker & \ding{55} & \ding{55} & \ding{55} & \ding{55} & \ding{51} & \ding{55} & \ding{51} & Produção \\
Rabbit Agents & \ding{55} & \ding{55} & \ding{51} & \ding{55} & \ding{55} & \ding{55} & \ding{51} & Produção \\
\hline
\textbf{MakerConnect (Proposto)} & \textbf{\ding{51}} & \textbf{\ding{51}} & \textbf{\ding{51}} & \textbf{\ding{51}} & \textbf{\ding{51}} & \textbf{\ding{51}} & \textbf{\ding{51}} & \textbf{MVP} \\
\hline
\end{tblr}
\vspace{2pt}
\caption*{\footnotesize Legenda: Prot.=Protótipo; Conc.=Conceitual. Critérios: \textit{Social}=presença de funcionalidades colaborativas da comunidade (p.ex., perfil, feed, comentários, upvotes, coautoria). \textit{Trace}=mecanismos explícitos de rastreabilidade técnica e linhagem de artefatos (p.ex., fork com parent\_id, histórico/versionamento, auditoria reprodutível).}
\label{tab:comparative}
\end{table}

\textbf{Análise:} a Tabela~\ref{tab:comparative} revela que nenhum trabalho antecedente combina todas as dimensões críticas. Especificamente:

\begin{itemize}
\item \textbf{RAG + IoT em conjunto:} Hangyu et al.~(2026) e Dong et al.~(2025) tratam ambos, mas nenhum integra mecanismos sociais de rastreabilidade nem orquestração complexa de múltiplos agentes.

\item \textbf{Orquestração de agentes:} Singh et al.~(2025) mapeia arquiteturas agentic RAG, mas foca em sistemas generalistas sem aplicação prática em IoT maker.

\item \textbf{Documentação técnica automatizada:} Dong et al.~(2025) avança ao gerar recomendações, mas como um assistente reativo de segurança, não como um pipeline pró-ativo para exportação de documentação padronizada (PDF).

\item \textbf{Rastreabilidade social (fork/lineage):} nenhum trabalho literário ou plataforma (incluindo Manual Maker) integra mecanismos sociais avançados (upvotes, forks com history) como camada de governança assistida.

\item \textbf{LGPD + compliance:} Hangyu et al.~(2026) aborda privacidade em ambientes distribuídos, mas não foca na proteção de dados de makers (Pessoa Física) no contexto de compartilhamento público de projetos.
\end{itemize}

\textbf{Posicionamento da MakerConnect:} O projeto preenche essa lacuna ao propor uma aplicação integrada que combina: (i)~RAG + IoT para documentação assistida; (ii)~orquestração n8n de agentes Generativos; (iii)~mecanismos sociais (fork lineage, upvotes, coautoria) para rastreabilidade; (iv)~conformidade LGPD com anonimização; (v)~implementação funcional MVP com validação em comunidade maker real.

% ============================================================
\section{Aspectos Éticos, LGPD e Qualidade}
% ============================================================

A conformidade com a LGPD é garantida por três camadas: (i)~\textbf{anonimização} de PII antes de qualquer processamento pela IA; (ii)~\textbf{consentimento explícito} para compartilhamento de projetos; (iii)~\textbf{criptografia} de dados em trânsito e em repouso. Os dados de componentes utilizados no RAG são estritamente técnicos, de domínio público e sem identificação pessoal.

Para atender aos critérios de qualidade desejáveis, a plataforma prevê: métricas de desempenho documental (cobertura de requisitos, tempo médio de geração, índice de reprodutibilidade de \textit{forks}); validação por \textit{holdout} sobre a relevância dos documentos recuperados pelo RAG; e infraestrutura de IA em nuvem (Hostinger, com possibilidade de migração para AWS/Azure).

% ============================================================
\section{Conclusão e Próximas Etapas}
% ============================================================

A MakerConnect demonstra como IA Generativa com RAG pode transformar uma rede social maker em ferramenta ativa de governança técnica, reduzindo a documentação fantasma em projetos IoT. O pipeline funcional atende integralmente aos requisitos mínimos de habilitação: abordagem fundamentada (RAG/Transformer), propósito prático integrado à aplicação, base de dados real de componentes, pipeline completo (extração → pré-proc. → modelo → pós-proc.) e conformidade LGPD. As próximas etapas incluem validação empírica com grupos de makers reais e publicação dos resultados.

% --- DECLARAÇÕES OBRIGATÓRIAS ---
\begin{declarations}
\begin{contributions}
VF foi responsável pela arquitetura do pipeline de IA, configuração do ambiente n8n, modelagem do banco de dados vetorial e desenvolvimento full-stack.
\end{contributions}
\begin{interests}
O autor declara não haver conflitos de interesse.
\end{interests}
\begin{materials}
O repositório ``Froesv85/Master-Labs/'' no GitHub conterá o código-fonte e as definições de workflow do n8n.
\end{materials}
\end{declarations}

\bibliographystyle{apalike-sol}
\bibliography{refs}

\end{document}